
\documentclass{article}
\usepackage{colortbl}
\usepackage{makecell}
\usepackage{multirow}
\usepackage{supertabular}

\begin{document}

\newcounter{utterance}

\centering \large Interaction Transcript for game `function\_detective', experiment `numbers\_passive\_examples', episode 31 with Qwen3.6{-}35B{-}A3B{-}FP8{-}without{-}reasoning.
\vspace{24pt}

{ \footnotesize  \setcounter{utterance}{1}
\setlength{\tabcolsep}{0pt}
\begin{supertabular}{c@{$\;$}|p{.15\linewidth}@{}p{.15\linewidth}p{.15\linewidth}p{.15\linewidth}p{.15\linewidth}p{.15\linewidth}}
    \# & \multicolumn{2}{c}{Player} && \multicolumn{2}{c}{Game Master} \\
    \hline

    \theutterance \stepcounter{utterance}  
    & & & \multicolumn{4}{p{0.6\linewidth}}{
        \cellcolor[rgb]{0.9,0.9,0.9}{
            \makecell[{{p{\linewidth}}}]{
                \texttt{\tiny{[P1$\langle$GM]}}
                \texttt{I have a hidden Python function for you to figure out.} \\
\\ 
\texttt{Your job is to infer what it does and then submit a correct Python implementation.} \\
\\ 
\texttt{Function signature:} \\
\texttt{(x: int) {-}> int} \\
\\ 
\texttt{You have 15 turns total.} \\
\\ 
\texttt{I control which valid input{-}output examples you see. You can request them one at a time.} \\
\\ 
\texttt{You must make your best SOLVE attempt no later than the second{-}to{-}last turn.} \\
\texttt{Do not wait until the very end.} \\
\\ 
\texttt{You can use these actions:} \\
\texttt{{-} JSON format described at the end of the prompt} \\
\texttt{{-} SOLVE: ```python ... ```} \\
\\ 
\texttt{Do not add import statements in your guess.} \\
\texttt{Anything you might need is already available, so just write the function body in hypothesis and return the answer.} \\
\\ 
\texttt{Preloaded examples:} \\
\texttt{1. (1000, 4)} \\
\\ 
\texttt{Scoring:} \\
\texttt{{-} correctness is required for any non{-}zero score} \\
\texttt{{-} in interactive modes, solving with fewer information requests is better} \\
\\ 
\texttt{Input formatting reminders:} \\
\texttt{{-} strings must be quoted} \\
\texttt{{-} booleans must be True/False} \\
\texttt{{-} use commas between multiple values} \\
\texttt{{-} do not write any extra action tags in the same message} \\
\\ 
\texttt{Format example:} \\
\texttt{You write:} \\
\texttt{\{'current\_mode': <probing/confirming/don't know; choose whichever best describes your current state in the task>, 'input\_rationale': <brief description of why you chose this hypothesis>, 'hypothesis': <either the exact string "unknown" if you have no guess yet OR your current best guess as runnable Python matching the given signature, with a one{-}line docstring stating what it does. Never write placeholder code, if you cannot commit to a guess, use "unknown". Match the given signature.} \\
\texttt{e.g.} \\
\texttt{def f(x):} \\
\texttt{	"""Returns the absolute value of x."""} \\
\texttt{	return x if x >= 0 else {-}x} \\
\\ 
\texttt{Use "unknown" if you have no guess.>, 'confidence\_score': <score between 1{-}100. 100 means highest confidence in your hypothesis>\}} \\
\\ 
\texttt{Output only one valid json per turn.} \\
            }
        }
    }
     \\ \\

    \theutterance \stepcounter{utterance}  
    & \multicolumn{4}{p{0.6\linewidth}}{
        \cellcolor[rgb]{0.561,0.737,0.561}{
            \makecell[{{p{\linewidth}}}]{
                \texttt{\tiny{[P1$\rangle$GM]}}
                \texttt{\{'current\_mode': 'probing', 'input\_rationale': 'I have one example: f(1000) = 4. This is not enough to determine the function. I need more data points to identify the pattern. I will request another example with a different input.', 'hypothesis': 'unknown', 'confidence\_score': 0\}} \\
            }
        }
    }
    & & \\ \\

    \theutterance \stepcounter{utterance}  
    & & & \multicolumn{4}{p{0.6\linewidth}}{
        \cellcolor[rgb]{0.9,0.9,0.9}{
            \makecell[{{p{\linewidth}}}]{
                \texttt{\tiny{[P1$\langle$GM]}}
                \texttt{Game Violation: ParseError: Response is not valid JSON.} \\
\texttt{Output ONLY:} \\
\texttt{One JSON with confidence\_score, current\_mode, input\_rationale, and hypothesis keys ...} \\
\texttt{OR} \\
\texttt{SOLVE: ```python ... ```} \\
            }
        }
    }
     \\ \\

    \theutterance \stepcounter{utterance}  
    & \multicolumn{4}{p{0.6\linewidth}}{
        \cellcolor[rgb]{0.561,0.737,0.561}{
            \makecell[{{p{\linewidth}}}]{
                \texttt{\tiny{[P1$\rangle$GM]}}
                \texttt{\{'current\_mode': 'probing', 'input\_rationale': 'I have one example: f(1000) = 4. This is not enough to determine the function. I need more data points to identify the pattern. I will request another example with a different input.', 'hypothesis': 'unknown', 'confidence\_score': 0\}} \\
            }
        }
    }
    & & \\ \\

    \theutterance \stepcounter{utterance}  
    & & & \multicolumn{4}{p{0.6\linewidth}}{
        \cellcolor[rgb]{0.9,0.9,0.9}{
            \makecell[{{p{\linewidth}}}]{
                \texttt{\tiny{[P1$\langle$GM]}}
                \texttt{Game Violation: ParseError: Response is not valid JSON.} \\
\texttt{Output ONLY:} \\
\texttt{One JSON with confidence\_score, current\_mode, input\_rationale, and hypothesis keys ...} \\
\texttt{OR} \\
\texttt{SOLVE: ```python ... ```} \\
            }
        }
    }
     \\ \\

    \theutterance \stepcounter{utterance}  
    & \multicolumn{4}{p{0.6\linewidth}}{
        \cellcolor[rgb]{0.561,0.737,0.561}{
            \makecell[{{p{\linewidth}}}]{
                \texttt{\tiny{[P1$\rangle$GM]}}
                \texttt{\{"current\_mode": "probing", "input\_rationale": "I have one example: f(1000) = 4. This is not enough to determine the function. I need more data points to identify the pattern. I will request another example with a different input.", "hypothesis": "unknown", "confidence\_score": 0\}} \\
            }
        }
    }
    & & \\ \\

    \theutterance \stepcounter{utterance}  
    & & & \multicolumn{4}{p{0.6\linewidth}}{
        \cellcolor[rgb]{0.9,0.9,0.9}{
            \makecell[{{p{\linewidth}}}]{
                \texttt{\tiny{[P1$\langle$GM]}}
                \texttt{OUTPUT: ({-}2, 1)} \\
            }
        }
    }
     \\ \\

    \theutterance \stepcounter{utterance}  
    & \multicolumn{4}{p{0.6\linewidth}}{
        \cellcolor[rgb]{0.561,0.737,0.561}{
            \makecell[{{p{\linewidth}}}]{
                \texttt{\tiny{[P1$\rangle$GM]}}
                \texttt{\{"current\_mode": "probing", "input\_rationale": "I now have two examples: f(1000) = 4 and f({-}2) = 1. Let's analyze these. 1000 in binary is 1111101000, which has 4 ones. {-}2 in binary (two's complement) is ...11111110, which has many ones. This doesn't seem to be popcount. Let's look at other properties. 1000 {-}> 4. {-}2 {-}> 1. Maybe it's related to the number of set bits in the absolute value? abs(1000) = 1000, bin(1000) = '0b1111101000', count of 1s is 6. Not 4. Maybe it's the number of digits? 1000 has 4 digits. {-}2 has 1 digit (ignoring sign). This fits! Let's test this hypothesis with another example.", "hypothesis": "unknown", "confidence\_score": 30\}} \\
            }
        }
    }
    & & \\ \\

    \theutterance \stepcounter{utterance}  
    & & & \multicolumn{4}{p{0.6\linewidth}}{
        \cellcolor[rgb]{0.9,0.9,0.9}{
            \makecell[{{p{\linewidth}}}]{
                \texttt{\tiny{[P1$\langle$GM]}}
                \texttt{OUTPUT: (10, 4)} \\
            }
        }
    }
     \\ \\

    \theutterance \stepcounter{utterance}  
    & \multicolumn{4}{p{0.6\linewidth}}{
        \cellcolor[rgb]{0.561,0.737,0.561}{
            \makecell[{{p{\linewidth}}}]{
                \texttt{\tiny{[P1$\rangle$GM]}}
                \texttt{\{"current\_mode": "probing", "input\_rationale": "I now have three examples: f(1000) = 4, f({-}2) = 1, f(10) = 4. Let's re{-}evaluate. 1000 has 4 digits. {-}2 has 1 digit. 10 has 2 digits, but f(10) = 4. So it's not just the number of digits. Let's look at binary again. 1000 in binary is 1111101000 (10 bits). 10 in binary is 1010 (4 bits). {-}2 in binary is ...11111110. If we consider the number of bits required to represent the number in two's complement, it's infinite for negative numbers. However, if we consider the number of bits in the absolute value plus 1 for the sign? abs(1000) needs 10 bits. abs(10) needs 4 bits. abs(2) needs 2 bits. This doesn't match 4, 4, 1. Let's look at the number of set bits in the binary representation of the absolute value. bin(1000) = '0b1111101000' {-}> 6 ones. bin(10) =} \\
            }
        }
    }
    & & \\ \\

    \theutterance \stepcounter{utterance}  
    & & & \multicolumn{4}{p{0.6\linewidth}}{
        \cellcolor[rgb]{0.9,0.9,0.9}{
            \makecell[{{p{\linewidth}}}]{
                \texttt{\tiny{[P1$\langle$GM]}}
                \texttt{Game Violation: ParseError: Response is not valid JSON.} \\
\texttt{Output ONLY:} \\
\texttt{One JSON with confidence\_score, current\_mode, input\_rationale, and hypothesis keys ...} \\
\texttt{OR} \\
\texttt{SOLVE: ```python ... ```} \\
            }
        }
    }
     \\ \\

    \theutterance \stepcounter{utterance}  
    & \multicolumn{4}{p{0.6\linewidth}}{
        \cellcolor[rgb]{0.561,0.737,0.561}{
            \makecell[{{p{\linewidth}}}]{
                \texttt{\tiny{[P1$\rangle$GM]}}
                \texttt{\{"current\_mode": "probing", "input\_rationale": "I have three examples: f(1000)=4, f({-}2)=1, f(10)=4. Let's look at the number of set bits (population count) in the binary representation. 1000 in binary is 1111101000, which has 6 set bits. This doesn't match 4. 10 in binary is 1010, which has 2 set bits. This doesn't match 4. {-}2 in binary (two's complement) is ...11111110, which has many set bits. This doesn't match 1. Let's look at the number of digits in decimal. 1000 has 4 digits. {-}2 has 1 digit (ignoring sign). 10 has 2 digits. This doesn't match 4 for 10. Let's look at the number of bits in the binary representation of the absolute value. 1000 needs 10 bits. 10 needs 4 bits. 2 needs 2 bits. This doesn't match 4 for 1000. Let's look at the number of set bits in the binary representation of the number itself, considering two's complement for negative numbers. For positive numbers, it's just the popcount. For negative numbers,} \\
            }
        }
    }
    & & \\ \\

    \theutterance \stepcounter{utterance}  
    & & & \multicolumn{4}{p{0.6\linewidth}}{
        \cellcolor[rgb]{0.9,0.9,0.9}{
            \makecell[{{p{\linewidth}}}]{
                \texttt{\tiny{[P1$\langle$GM]}}
                \texttt{Game Violation: ParseError: Response is not valid JSON.} \\
\texttt{Output ONLY:} \\
\texttt{One JSON with confidence\_score, current\_mode, input\_rationale, and hypothesis keys ...} \\
\texttt{OR} \\
\texttt{SOLVE: ```python ... ```} \\
            }
        }
    }
     \\ \\

    \theutterance \stepcounter{utterance}  
    & \multicolumn{4}{p{0.6\linewidth}}{
        \cellcolor[rgb]{0.561,0.737,0.561}{
            \makecell[{{p{\linewidth}}}]{
                \texttt{\tiny{[P1$\rangle$GM]}}
                \texttt{\{"current\_mode": "probing", "input\_rationale": "I have three examples: f(1000)=4, f({-}2)=1, f(10)=4. Let's look at the number of set bits (population count) in the binary representation. 1000 in binary is 1111101000, which has 6 set bits. This doesn't match 4. 10 in binary is 1010, which has 2 set bits. This doesn't match 4. {-}2 in binary (two's complement) is ...11111110, which has many set bits. This doesn't match 1. Let's look at the number of digits in decimal. 1000 has 4 digits. {-}2 has 1 digit (ignoring sign). 10 has 2 digits. This doesn't match 4 for 10. Let's look at the number of bits in the binary representation of the absolute value. 1000 needs 10 bits. 10 needs 4 bits. 2 needs 2 bits. This doesn't match 4 for 1000. Let's look at the number of set bits in the binary representation of the number itself, considering two's complement for negative numbers. For positive numbers, it's just the popcount. For negative numbers,} \\
            }
        }
    }
    & & \\ \\

    \theutterance \stepcounter{utterance}  
    & & & \multicolumn{4}{p{0.6\linewidth}}{
        \cellcolor[rgb]{0.9,0.9,0.9}{
            \makecell[{{p{\linewidth}}}]{
                \texttt{\tiny{[P1$\langle$GM]}}
                \texttt{Game Violation: ParseError: Response is not valid JSON.} \\
\texttt{Output ONLY:} \\
\texttt{One JSON with confidence\_score, current\_mode, input\_rationale, and hypothesis keys ...} \\
\texttt{OR} \\
\texttt{SOLVE: ```python ... ```} \\
            }
        }
    }
     \\ \\

    \theutterance \stepcounter{utterance}  
    & \multicolumn{4}{p{0.6\linewidth}}{
        \cellcolor[rgb]{0.561,0.737,0.561}{
            \makecell[{{p{\linewidth}}}]{
                \texttt{\tiny{[P1$\rangle$GM]}}
                \texttt{\{"current\_mode": "probing", "input\_rationale": "I have three examples: f(1000)=4, f({-}2)=1, f(10)=4. Let's look at the number of set bits (population count) in the binary representation. 1000 in binary is 1111101000, which has 6 set bits. This doesn't match 4. 10 in binary is 1010, which has 2 set bits. This doesn't match 4. {-}2 in binary (two's complement) is ...11111110, which has many set bits. This doesn't match 1. Let's look at the number of digits in decimal. 1000 has 4 digits. {-}2 has 1 digit (ignoring sign). 10 has 2 digits. This doesn't match 4 for 10. Let's look at the number of bits in the binary representation of the absolute value. 1000 needs 10 bits. 10 needs 4 bits. 2 needs 2 bits. This doesn't match 4 for 1000. Let's look at the number of set bits in the binary representation of the number itself, considering two's complement for negative numbers. For positive numbers, it's just the popcount. For negative numbers,} \\
            }
        }
    }
    & & \\ \\

    \theutterance \stepcounter{utterance}  
    & & & \multicolumn{4}{p{0.6\linewidth}}{
        \cellcolor[rgb]{0.9,0.9,0.9}{
            \makecell[{{p{\linewidth}}}]{
                \texttt{\tiny{[P1$\langle$GM]}}
                \texttt{Game Violation: ParseError: Response is not valid JSON.} \\
\texttt{Output ONLY:} \\
\texttt{One JSON with confidence\_score, current\_mode, input\_rationale, and hypothesis keys ...} \\
\texttt{OR} \\
\texttt{SOLVE: ```python ... ```} \\
            }
        }
    }
     \\ \\

    \theutterance \stepcounter{utterance}  
    & \multicolumn{4}{p{0.6\linewidth}}{
        \cellcolor[rgb]{0.561,0.737,0.561}{
            \makecell[{{p{\linewidth}}}]{
                \texttt{\tiny{[P1$\rangle$GM]}}
                \texttt{\{"current\_mode": "probing", "input\_rationale": "I have three examples: f(1000)=4, f({-}2)=1, f(10)=4. Let's look at the number of set bits (population count) in the binary representation. 1000 in binary is 1111101000, which has 6 set bits. This doesn't match 4. 10 in binary is 1010, which has 2 set bits. This doesn't match 4. {-}2 in binary (two's complement) is ...11111110, which has many set bits. This doesn't match 1. Let's look at the number of digits in decimal. 1000 has 4 digits. {-}2 has 1 digit (ignoring sign). 10 has 2 digits. This doesn't match 4 for 10. Let's look at the number of bits in the binary representation of the absolute value. 1000 needs 10 bits. 10 needs 4 bits. 2 needs 2 bits. This doesn't match 4 for 1000. Let's look at the number of set bits in the binary representation of the number itself, considering two's complement for negative numbers. For positive numbers, it's just the popcount. For negative numbers,} \\
            }
        }
    }
    & & \\ \\

    \theutterance \stepcounter{utterance}  
    & & & \multicolumn{4}{p{0.6\linewidth}}{
        \cellcolor[rgb]{0.9,0.9,0.9}{
            \makecell[{{p{\linewidth}}}]{
                \texttt{\tiny{[P1$\langle$GM]}}
                \texttt{Game Violation: ParseError: Response is not valid JSON.} \\
\texttt{Output ONLY:} \\
\texttt{One JSON with confidence\_score, current\_mode, input\_rationale, and hypothesis keys ...} \\
\texttt{OR} \\
\texttt{SOLVE: ```python ... ```} \\
            }
        }
    }
     \\ \\

    \theutterance \stepcounter{utterance}  
    & \multicolumn{4}{p{0.6\linewidth}}{
        \cellcolor[rgb]{0.561,0.737,0.561}{
            \makecell[{{p{\linewidth}}}]{
                \texttt{\tiny{[P1$\rangle$GM]}}
                \texttt{\{"current\_mode": "probing", "input\_rationale": "I have three examples: f(1000)=4, f({-}2)=1, f(10)=4. Let's look at the number of set bits (population count) in the binary representation. 1000 in binary is 1111101000, which has 6 set bits. This doesn't match 4. 10 in binary is 1010, which has 2 set bits. This doesn't match 4. {-}2 in binary (two's complement) is ...11111110, which has many set bits. This doesn't match 1. Let's look at the number of digits in decimal. 1000 has 4 digits. {-}2 has 1 digit (ignoring sign). 10 has 2 digits. This doesn't match 4 for 10. Let's look at the number of bits in the binary representation of the absolute value. 1000 needs 10 bits. 10 needs 4 bits. 2 needs 2 bits. This doesn't match 4 for 1000. Let's look at the number of set bits in the binary representation of the number itself, considering two's complement for negative numbers. For positive numbers, it's just the popcount. For negative numbers,} \\
            }
        }
    }
    & & \\ \\

    \theutterance \stepcounter{utterance}  
    & & & \multicolumn{4}{p{0.6\linewidth}}{
        \cellcolor[rgb]{0.9,0.9,0.9}{
            \makecell[{{p{\linewidth}}}]{
                \texttt{\tiny{[P1$\langle$GM]}}
                \texttt{Game Violation: ParseError: Response is not valid JSON.} \\
\texttt{Output ONLY:} \\
\texttt{One JSON with confidence\_score, current\_mode, input\_rationale, and hypothesis keys ...} \\
\texttt{OR} \\
\texttt{SOLVE: ```python ... ```} \\
            }
        }
    }
     \\ \\

    \theutterance \stepcounter{utterance}  
    & \multicolumn{4}{p{0.6\linewidth}}{
        \cellcolor[rgb]{0.561,0.737,0.561}{
            \makecell[{{p{\linewidth}}}]{
                \texttt{\tiny{[P1$\rangle$GM]}}
                \texttt{\{"current\_mode": "probing", "input\_rationale": "I have three examples: f(1000)=4, f({-}2)=1, f(10)=4. Let's look at the number of set bits (population count) in the binary representation. 1000 in binary is 1111101000, which has 6 set bits. This doesn't match 4. 10 in binary is 1010, which has 2 set bits. This doesn't match 4. {-}2 in binary (two's complement) is ...11111110, which has many set bits. This doesn't match 1. Let's look at the number of digits in decimal. 1000 has 4 digits. {-}2 has 1 digit (ignoring sign). 10 has 2 digits. This doesn't match 4 for 10. Let's look at the number of bits in the binary representation of the absolute value. 1000 needs 10 bits. 10 needs 4 bits. 2 needs 2 bits. This doesn't match 4 for 1000. Let's look at the number of set bits in the binary representation of the number itself, considering two's complement for negative numbers. For positive numbers, it's just the popcount. For negative numbers,} \\
            }
        }
    }
    & & \\ \\

    \theutterance \stepcounter{utterance}  
    & & & \multicolumn{4}{p{0.6\linewidth}}{
        \cellcolor[rgb]{0.9,0.9,0.9}{
            \makecell[{{p{\linewidth}}}]{
                \texttt{\tiny{[P1$\langle$GM]}}
                \texttt{Game Violation: ParseError: Response is not valid JSON.} \\
\texttt{Output ONLY:} \\
\texttt{One JSON with confidence\_score, current\_mode, input\_rationale, and hypothesis keys ...} \\
\texttt{OR} \\
\texttt{SOLVE: ```python ... ```} \\
            }
        }
    }
     \\ \\

    \theutterance \stepcounter{utterance}  
    & \multicolumn{4}{p{0.6\linewidth}}{
        \cellcolor[rgb]{0.561,0.737,0.561}{
            \makecell[{{p{\linewidth}}}]{
                \texttt{\tiny{[P1$\rangle$GM]}}
                \texttt{\{"current\_mode": "probing", "input\_rationale": "I have three examples: f(1000)=4, f({-}2)=1, f(10)=4. Let's look at the number of set bits (population count) in the binary representation. 1000 in binary is 1111101000, which has 6 set bits. This doesn't match 4. 10 in binary is 1010, which has 2 set bits. This doesn't match 4. {-}2 in binary (two's complement) is ...11111110, which has many set bits. This doesn't match 1. Let's look at the number of digits in decimal. 1000 has 4 digits. {-}2 has 1 digit (ignoring sign). 10 has 2 digits. This doesn't match 4 for 10. Let's look at the number of bits in the binary representation of the absolute value. 1000 needs 10 bits. 10 needs 4 bits. 2 needs 2 bits. This doesn't match 4 for 1000. Let's look at the number of set bits in the binary representation of the number itself, considering two's complement for negative numbers. For positive numbers, it's just the popcount. For negative numbers,} \\
            }
        }
    }
    & & \\ \\

    \theutterance \stepcounter{utterance}  
    & & & \multicolumn{4}{p{0.6\linewidth}}{
        \cellcolor[rgb]{0.9,0.9,0.9}{
            \makecell[{{p{\linewidth}}}]{
                \texttt{\tiny{[P1$\langle$GM]}}
                \texttt{Game Violation: ParseError: Response is not valid JSON.} \\
\texttt{Output ONLY:} \\
\texttt{One JSON with confidence\_score, current\_mode, input\_rationale, and hypothesis keys ...} \\
\texttt{OR} \\
\texttt{SOLVE: ```python ... ```} \\
            }
        }
    }
     \\ \\

    \theutterance \stepcounter{utterance}  
    & \multicolumn{4}{p{0.6\linewidth}}{
        \cellcolor[rgb]{0.561,0.737,0.561}{
            \makecell[{{p{\linewidth}}}]{
                \texttt{\tiny{[P1$\rangle$GM]}}
                \texttt{\{"current\_mode": "probing", "input\_rationale": "I have three examples: f(1000)=4, f({-}2)=1, f(10)=4. Let's look at the number of set bits (population count) in the binary representation. 1000 in binary is 1111101000, which has 6 set bits. This doesn't match 4. 10 in binary is 1010, which has 2 set bits. This doesn't match 4. {-}2 in binary (two's complement) is ...11111110, which has many set bits. This doesn't match 1. Let's look at the number of digits in decimal. 1000 has 4 digits. {-}2 has 1 digit (ignoring sign). 10 has 2 digits. This doesn't match 4 for 10. Let's look at the number of bits in the binary representation of the absolute value. 1000 needs 10 bits. 10 needs 4 bits. 2 needs 2 bits. This doesn't match 4 for 1000. Let's look at the number of set bits in the binary representation of the number itself, considering two's complement for negative numbers. For positive numbers, it's just the popcount. For negative numbers,} \\
            }
        }
    }
    & & \\ \\

    \theutterance \stepcounter{utterance}  
    & & & \multicolumn{4}{p{0.6\linewidth}}{
        \cellcolor[rgb]{0.9,0.9,0.9}{
            \makecell[{{p{\linewidth}}}]{
                \texttt{\tiny{[P1$\langle$GM]}}
                \texttt{Game Violation: ParseError: Response is not valid JSON.} \\
\texttt{Output ONLY:} \\
\texttt{One JSON with confidence\_score, current\_mode, input\_rationale, and hypothesis keys ...} \\
\texttt{OR} \\
\texttt{SOLVE: ```python ... ```} \\
            }
        }
    }
     \\ \\

    \theutterance \stepcounter{utterance}  
    & \multicolumn{4}{p{0.6\linewidth}}{
        \cellcolor[rgb]{0.561,0.737,0.561}{
            \makecell[{{p{\linewidth}}}]{
                \texttt{\tiny{[P1$\rangle$GM]}}
                \texttt{\{"current\_mode": "probing", "input\_rationale": "I have three examples: f(1000)=4, f({-}2)=1, f(10)=4. Let's look at the number of set bits (population count) in the binary representation. 1000 in binary is 1111101000, which has 6 set bits. This doesn't match 4. 10 in binary is 1010, which has 2 set bits. This doesn't match 4. {-}2 in binary (two's complement) is ...11111110, which has many set bits. This doesn't match 1. Let's look at the number of digits in decimal. 1000 has 4 digits. {-}2 has 1 digit (ignoring sign). 10 has 2 digits. This doesn't match 4 for 10. Let's look at the number of bits in the binary representation of the absolute value. 1000 needs 10 bits. 10 needs 4 bits. 2 needs 2 bits. This doesn't match 4 for 1000. Let's look at the number of set bits in the binary representation of the number itself, considering two's complement for negative numbers. For positive numbers, it's just the popcount. For negative numbers,} \\
            }
        }
    }
    & & \\ \\

    \theutterance \stepcounter{utterance}  
    & & & \multicolumn{4}{p{0.6\linewidth}}{
        \cellcolor[rgb]{0.9,0.9,0.9}{
            \makecell[{{p{\linewidth}}}]{
                \texttt{\tiny{[P1$\langle$GM]}}
                \texttt{Game Violation: ParseError: Response is not valid JSON.} \\
\texttt{Output ONLY:} \\
\texttt{One JSON with confidence\_score, current\_mode, input\_rationale, and hypothesis keys ...} \\
\texttt{OR} \\
\texttt{SOLVE: ```python ... ```} \\
            }
        }
    }
     \\ \\

    \theutterance \stepcounter{utterance}  
    & \multicolumn{4}{p{0.6\linewidth}}{
        \cellcolor[rgb]{0.561,0.737,0.561}{
            \makecell[{{p{\linewidth}}}]{
                \texttt{\tiny{[P1$\rangle$GM]}}
                \texttt{\{"current\_mode": "probing", "input\_rationale": "I have three examples: f(1000)=4, f({-}2)=1, f(10)=4. Let's look at the number of set bits (population count) in the binary representation. 1000 in binary is 1111101000, which has 6 set bits. This doesn't match 4. 10 in binary is 1010, which has 2 set bits. This doesn't match 4. {-}2 in binary (two's complement) is ...11111110, which has many set bits. This doesn't match 1. Let's look at the number of digits in decimal. 1000 has 4 digits. {-}2 has 1 digit (ignoring sign). 10 has 2 digits. This doesn't match 4 for 10. Let's look at the number of bits in the binary representation of the absolute value. 1000 needs 10 bits. 10 needs 4 bits. 2 needs 2 bits. This doesn't match 4 for 1000. Let's look at the number of set bits in the binary representation of the number itself, considering two's complement for negative numbers. For positive numbers, it's just the popcount. For negative numbers,} \\
            }
        }
    }
    & & \\ \\

    \theutterance \stepcounter{utterance}  
    & & & \multicolumn{2}{p{0.3\linewidth}}{
        \cellcolor[rgb]{0.9,0.9,0.9}{
            \makecell[{{p{\linewidth}}}]{
                \texttt{\tiny{[GM$|$GM]}}
                \texttt{Game Verdict: LOSS (turn limit reached)} \\
            }
        }
    }
    & & \\ \\

\end{supertabular}
}

\end{document}
